\scanpage{26}
la $=$ propriété de positivité, savoir
\[
\mathrm{Tr}\,u^\psi u\deleted{\unsure{=}}\quad(=\mathrm{Tr}\,(a^{-1}ua)^\varphi(a^{-1}ua))\geq0
\]
égalité seulement si $u=0$. Donc tt pt de
l'intérieur du cône \deleted{de} $\Phi^+$ engendre
(via autom de $M$) le $=$ cône $\Phi^+$.

C'est la donnée du cône $\Phi^+$ qui,
intrinsèquement, constitue une structure
supplémentaire importante.
% slanted
\margin{dégager la structure abstraite : anneau $A$ \inserted{de type fini sur $Q$}, bimodule $\Phi$ sur $A$ libre à gauche et à droite, Trace sur $A$ ; notion de \uline{bon} cône ds $\Phi$}

On notera que, sous cette forme abstraite,
on ne peut (parmi les cônes possibles,
correspondant aux divers choix de $\varphi$)
en distinguer de meilleurs que d'autres.
Ainsi, si $\varphi$ satisfait la condition de
positivité de départ, de $=$ $-\varphi$, qui
définira le cône $-\Phi^+$

Noter aussi : \deleted{Tout cône convexe fermé
qui contient} si $v,w\in A^+$,
alors $\mathrm{Tr}\,vw\geq0$ [et
en fait, $A^+$ est son \uline{propre} \uline{polaire} dans
$A^s$ \inserted{hermitiens} pour la forme $\mathrm{Tr}\,\unsure{uv|vw}$]. Or pour
% archive sheet number: 25
